\documentclass[11pt]{article}
\usepackage{mathblatt}
% Präambel des Setzers (werkzeuge/setzer-praeambel.tex, wortgleich); dazu \gteil (Muster T2U)
\makeatletter
\setlength{\mbpfbe}{0mm}% keine Punkte (pruefung.md § 5: nur beim Prüfstein)
\newlength{\gEin}\setlength{\gEin}{0pt}
\let\gSchrift\relax
\renewcommand{\pfkreuzzeile}[1]{\par\noindent\hangindent3.2em\hangafter1\hspace*{1.6em}$\square$\ #1\par}
\newcommand{\gkreuz}[1]{$\square$\ #1\hspace{2em}}
\newcommand{\gkreuzzeile}[2]{\par\noindent\hangindent2.2em\hangafter1\mbox{#1\,$\square$}\ #2\par}% Buchstabe, Option: eine je Zeile (neben Skizze)
% links, was man ansieht, rechts, was man tut; Spaltengrenze überall an derselben Stelle (gemeinsam.md „Satz“)
% v0.17: Fußeintrag „Nr: Wert“ – kurz in einer Box (nie getrennt), lang umbrechbar mit der Nummer am Anfang
% (vorher \mbox: lange Lösungen liefen über den Rand, AS40 138 pt); robust, weil \fusshilfe in die .aux schreibt
\NewDocumentCommand{\gfuss}{m m}{\sbox0{#1:~#2}\ifdim\wd0<0.45\textwidth\mbox{#1:~#2}\else\mbox{#1:}\nobreak\ #2\fi}
\newsavebox{\gbox}\newlength{\gLinks}
% v0.10 (Lehrer 13.10.): keine starre Halbierung – die linke Spalte ist so breit wie die Figur (+ 6 mm), die
% Antwort steht direkt daneben; ist die Figur breiter als 0,62 Zeile, steht die Antwort darunter
\newcommand{\gzwei}[2]{\par\smallskip\noindent\sbox{\gbox}{#1}%
  \setlength{\gLinks}{\dimexpr\wd\gbox+6mm\relax}%
  \ifdim\gLinks>0.62\linewidth
    \usebox{\gbox}\par\smallskip\noindent\begin{minipage}{\linewidth}#2\end{minipage}%
  \else
    \begin{minipage}[t]{\gLinks}\vspace{0pt}\usebox{\gbox}\end{minipage}%
    \begin{minipage}[t]{\dimexpr\linewidth-\gLinks\relax}\vspace{0pt}\raggedright #2\end{minipage}%
  \fi\par}
% 16.10. (Lehrer, PA10 S. 1): Text links, Koordinatensystem rechts daneben; ist das System breiter als 0,58 Zeile,
% steht es wie bisher unter dem Text
% Passt das System knapp nicht mehr auf die Seite (bis 25 % zu hoch), wird es auf den Rest verkleinert statt
% die Seite halb leer zu lassen
\newlength{\gRest}\newlength{\gHoch}
\newcommand{\gtextks}[2]{\sbox{\gbox}{#1}%
  \setlength{\gHoch}{\dimexpr\ht\gbox+\dp\gbox\relax}%
  \ifdim\pagegoal<\maxdimen\setlength{\gRest}{\dimexpr\pagegoal-\pagetotal-4\baselineskip\relax}%
    \ifdim\gHoch>\gRest\ifdim\gRest>0.75\gHoch
      \sbox{\gbox}{\resizebox*{!}{\gRest}{\usebox{\gbox}}}\fi\fi\fi
  \ifdim\wd\gbox>0.58\linewidth #2\par\smallskip\noindent\gpassend{\usebox{\gbox}}\par
  \else\begin{minipage}[t]{\dimexpr\linewidth-\wd\gbox-5mm\relax}\raggedright #2\end{minipage}\hfill
    \begin{minipage}[t]{\wd\gbox}\vspace{0pt}\usebox{\gbox}\end{minipage}\par\fi}
% v0.9 (12.10., pruefung.md § 5 „Marke mit B, Bedingungen rechts“): drittes Argument von pfaufg = Bedingungen
% der Prüfung (v0.10: \pfohnetr = TR durchgestrichen, „FOR“), klein grau im rechten Rand auf der Grundlinie der Nummer, Zeilen mit \\
\newcommand{\gbedrand}[2]{\ifblank{#2}{}{\rlap{\hspace*{\dimexpr#1+2mm\relax}%
  \parbox[t]{15mm}{\raggedright\scriptsize\leavevmode\color{mbgrau}\strut#2}}}}
% Teil eines Bündels (a, b …): Marke links im Rand, Bedingungen rechts im Rand, auf der Zeile des Teils
\newcommand{\gteilrand}[2]{\llap{\mbpfrandmarke{#1}\hspace*{\mbpfnr}}\gbedrand{\linewidth}{#2}}
% Teil eines Bündels: Buchstabe in eigener Spalte, Inhalt eingerückt (\gEin = Einzug, damit die Spaltengrenze von
% \gzwei an derselben Stelle bleibt)
\newenvironment{gbteil}[3]{\par\smallskip\noindent\gteilrand{#2}{#3}\makebox[\mbpfteil][l]{\textbf{#1}}%
  \begin{minipage}[t]{\dimexpr\linewidth-\mbpfteil\relax}\setlength{\gEin}{\mbpfteil}\raggedright\strut\ignorespaces}%
 {\end{minipage}\par\addvspace{6pt}}% v0.18 (PG80 Nr. 1): Abstand, damit die Linie von a) nicht an b) rutscht
% Grafik höchstens so breit wie die Spalte (Bündel nebeneinander, Koordinatensystem ohne Karo)
\newsavebox{\gpbox}
\newcommand{\gpassend}[1]{\sbox{\gpbox}{#1}\ifdim\wd\gpbox>\linewidth\typeout{GPASSEND: verkleinert \the\wd\gpbox\space auf \the\linewidth}\resizebox{\linewidth}{!}{\usebox{\gpbox}}\else\usebox{\gpbox}\fi}
% Skizzenfeld (v0.10, Lehrer 13.10.): Karo ab dem linken Rand über die Zeile, „Skizze“ klein grau oben links im Feld
% (wie \pfskizzenfeld der Vorlage); der rechte Rand bleibt den Bedingungen
\newcommand{\gskizze}[1]{\par\smallskip\noindent\begin{tikzpicture}
  \clip (0,0) rectangle ({\linewidth-1mm},{#1*10mm});
  \draw[black!16,line width=0.3pt,step=5mm] (0,0) grid ({\linewidth},{#1*10mm});
  \node[anchor=north west,inner sep=2pt,font=\scriptsize,text=mbgrau,fill=white] at (0,{#1*10mm}) {Skizze};
  \end{tikzpicture}\par}
\RenewDocumentEnvironment{pfaufg}{m m m}{%
  \begin{lrbox}{\mbaufgabebox}\begin{minipage}[t]{\linewidth}%
  \gSchrift\noindent\hspace*{\gEin}\mbpfrandmarke{#1}\gbedrand{\linewidth-\gEin}{#3}%
  \begin{minipage}[t]{\mbpfnr}\strut\textbf{#2}\end{minipage}%
  \begin{minipage}[t]{\dimexpr\linewidth-\gEin-\mbpfnr\relax}\raggedright\strut\ignorespaces}%
 {\end{minipage}%
  \end{minipage}\end{lrbox}%
  \setlength{\mbaufgabehoehe}{\dimexpr\ht\mbaufgabebox+\dp\mbaufgabebox\relax}%
  \ifdefempty{\gkopfnach}{\par\addvspace{5pt}\Needspace*{\mbaufgabehoehe}}{%
    % Stufenkopf, der mit dieser (ersten) Aufgabe zusammenbleibt: Platz für beide verlangen, dann den Kopf setzen
    % (12.10., Quadratische: ein Kopf stand allein unten, die Aufgabe mit Koordinatensystem auf der nächsten Seite)
    % v0.17 (AS40 „Terme deuten“ allein unten): hinter dem Kopf steht das Sprungziel (\special, ein Knoten, nach
    % dem die Seite brechen durfte) – darum \nobreak nach dem Kopf; Vorhalt 56 pt = Kopf mit Linie und Abständen
    \par\Needspace*{\dimexpr\mbaufgabehoehe+56pt\relax}\gkopfnach\par\nobreak\gdef\gkopfnach{}\addvspace{5pt}}%
  \noindent\usebox{\mbaufgabebox}\par\addvspace{5pt}}
% v0.18 (plan.md § 10, 15.10. l): Abschnittsbuchstabe A, B, C vor dem Namen im Kopf – als Vorsatz, damit der
% Name im TeX gleich bleibt (Bestellkarte, Sprungziele lesen ihn)
\def\gbstakt{}\newcommand{\gbstvor}[1]{\gdef\gbstakt{#1}}
\NewCommandCopy\gpfstufekopfalt\pfstufekopf
\RenewDocumentCommand{\pfstufekopf}{o m m}{\IfValueTF{#1}{\gpfstufekopfalt[#1]}{\gpfstufekopfalt}%
  {\ifx\gbstakt\empty\else\gbstakt\enspace\fi #2}{#3}\gdef\gbstakt{}}
% v0.18 (16.10. e): Begriff grau, ohne Nummer, ohne Rahmen über der ersten Aufgabe der Stufe
\newcommand{\gbegriff}[1]{\par\nobreak\noindent\hspace*{\mbpfnr}\parbox[t]{\dimexpr\linewidth-\mbpfnr\relax}{\raggedright\small\color{mbgrau}#1}\par\nobreak\addvspace{3pt}}
% v0.18 (16.10. g): Lösungsblatt – Kopf mit kräftigem Strich, Abschnitt klein grau als Zwischenzeile
\newcommand{\glkopf}[1]{\noindent{\large\bfseries Lösungen \textperiodcentered{} #1}\par\vspace{3pt}\hrule height 1.2pt\par\medskip}
\newcommand{\glzwischen}[1]{\par\addvspace{8pt}\Needspace*{4\baselineskip}\noindent{\footnotesize\color{mbgrau}#1}\par\nobreak\vspace{1pt}}
% v0.18 (15.10. o, w): Blattende – Kern ohne Rahmen, dann graue Zeile „Mehr: n Aufgaben“
\newcommand{\gkern}[1]{\par\addvspace{14pt}\Needspace*{6\baselineskip}\noindent{\bfseries Kern}\par\nobreak\smallskip
  {\small\raggedright #1\par}}
\newcommand{\gmehr}[1]{\par\addvspace{10pt}\noindent{\footnotesize\color{mbgrau}Mehr: #1 Aufgaben}\par}
% Stufenkopf vormerken: \gkopfvor{<Kopf>} – die nächste pfaufg setzt ihn, sobald beide zusammen Platz haben
\def\gkopfnach{}
\newcommand{\gkopfvor}[1]{\gdef\gkopfnach{#1}}
% Teilaufgabe mit Skizze links in gzwei: Buchstabe in eigener Spalte, darunter Text und Skizze
\newcommand{\gteil}[2]{\makebox[\mbpfteil][l]{\textbf{#1}}\parbox[t]{\dimexpr0.5\textwidth-\gEin-\mbpfnr-5mm-\mbpfteil\relax}{\raggedright #2}}
% Stufenkopf der Vorlage (mathblatt.sty 2026-10-12: Linie, Abstand, 13 pt, \pfjahre) ohne ihren Vorhalt von
% 14 Zeilen: der Setzer fordert vor dem Kopf 0,16 Seitenhöhe (Kopf + erste Aufgabe), mehr nicht
\setlength{\mbpfstufevorhalt}{0pt}
\makeatother
\newcommand{\karo}[1]{\par\smallskip\noindent\begin{tikzpicture}
  \clip (0,0) rectangle ({\linewidth-1mm},{#1*10mm});
  \draw[black!16,line width=0.3pt,step=5mm] (0,0) grid ({\linewidth},{#1*10mm});
  \end{tikzpicture}\par}
\newcommand{\kk}[1]{$\square$\,#1\hspace{0.9em}}
\newcommand{\linie}[1]{\,{\color{black!45}\rule[-1pt]{#1}{0.4pt}}\,}
\newcommand{\pkt}{\textperiodcentered{}}
% v0.14 (14.10.): Stufenkopf mit fünf Punkten (Option = Jahre für die Bestellkarte) oder „zuletzt 19“
\NewDocumentCommand{\gpunkte}{o m}{#2}
\newcommand{\gzuletzt}[1]{zuletzt #1}
% v0.7 (12.10.): Schülerfrage unter dem Stufenkopf; Ankreuzoption als Bild (Buchstabe und Kästchen darunter)
\newcommand{\gfrage}[1]{\par\nobreak\noindent\hspace*{\mbpfnr}{\small\itshape #1}\par\nobreak\addvspace{2pt}}
\newlength{\gbildbreite}\setlength{\gbildbreite}{0.32\linewidth}% Setzer setzt sie je Reihe (vier Bilder: 0,24)
\newcommand{\gbildoption}[3]{\begin{minipage}[t][#3][b]{\gbildbreite}\centering #2\par\smallskip\mbox{#1\,$\square$}\end{minipage}}
% Lineare (12.10.): Bild mit Buchstabe ohne Kästchen (Diagramm zuordnen); Aussage mit „wahr / falsch“ zum Ankreuzen
\newcommand{\gbildname}[3]{\begin{minipage}[t][#3][b]{\gbildbreite}\centering #2\par\smallskip\mbox{\textbf{#1}}\end{minipage}}
\newcommand{\gwahrfalsch}[2]{\par\noindent\hangindent2.2em\hangafter1\makebox[2.2em][l]{(#1)}#2\hfill\mbox{wahr $\square$\quad falsch $\square$}\par}
\makeatletter\setlength{\mbpfbe}{0mm}\makeatother
\begin{document}

\pfheftstil
\pflaufkopf[P10 \pkt{} mehr]{Parabel}{PA30-6}\pfseitevon
\pfheftkopf{Parabel}{mehr}
\pfstufenzeile{Scheitelpunkt angeben \pkt{} Scheitelpunktform aufstellen \pkt{} Parabel zeichnen \pkt{} Lage begründen}
% 2021-OS-B1d-m1 (Schwere I, rang 8)
\begin{pfaufg}{nach~P10\\’21~\pkt{}~1d}{1.}{}
Kreuze den Scheitelpunkt der gezeichneten Parabel an:
\par\smallskip\noindent\gpassend{\begin{ksys}[xmin=-6,xmax=1,ymin=-1,ymax=5,karo=0.5]\parabel{1}{-3}{1}{}\end{ksys}}\par
\par\smallskip\noindent\pfkreuz{}{$S(3\,|\,1)$}\pfkreuz{}{$S(-3\,|\,-1)$}\pfkreuz{}{$S(1\,|\,-3)$}\pfkreuz{}{$S(-3\,|\,1)$.}
\fusshilfe{\gfuss{1}{S($-$3|1)}}
\end{pfaufg}
% 2020-OS-K3a-m1 (Schwere I, rang 9)
\begin{pfaufg}{nach~P10\\’20~\pkt{}~3a}{2.}{}
Die Parabel zu $f$ mit $y = (x+3)^2 - 5$ ist im Koordinatensystem gezeichnet. $S$ ist der Scheitelpunkt des Graphen von $f$. Gib seine Koordinaten an.
\par\smallskip\noindent\gpassend{\begin{ksys}[xmin=-7,xmax=7,ymin=-6,ymax=5,karo=0.5]\parabel{1}{-3}{-5}{}\node[fill=white,inner sep=1pt,font=\small] at (-6.5,3) {$f$};\end{ksys}}\par
\karo{2}
\fusshilfe{\gfuss{2}{S($-$3|$-$5)}}
\end{pfaufg}
% 2018-OS-K5b-m1 (Schwere I, rang 12)
\begin{pfaufg}{nach~P10\\’18~\pkt{}~5b}{3.}{}
Im Koordinatensystem ist die Parabel $p$ mit $p(x) = x^2 - 6x + 6$ gezeichnet. Gib den Scheitelpunkt von $p$ an: $S(\ \ |\ \ )$.
\par\smallskip\noindent\gpassend{\begin{ksys}[xmin=-3,xmax=7,ymin=-4,ymax=7,karo=0.5]\parabel{1}{3}{-3}{p}\end{ksys}}\par
\karo{2}
\fusshilfe{\gfuss{3}{S(3|$-$3)}}
\end{pfaufg}
% 2017-OS-K5c-m1 (Schwere I, rang 13)
\begin{pfaufg}{nach~P10\\’17~\pkt{}~5c}{4.}{}
Die Parabel $p$ hat die Gleichung $y = (x+5)^2 - 3$ und ist im Koordinatensystem gezeichnet. Gib den Scheitelpunkt der Parabel $p$ an: $S(\ \ |\ \ )$.
\par\smallskip\noindent\gpassend{\begin{ksys}[xmin=-9,xmax=1,ymin=-4,ymax=4,karo=0.5]\parabel{1}{-5}{-3}{p}\end{ksys}}\par
\karo{2}
\fusshilfe{\gfuss{4}{S($-$5|$-$3)}}
\end{pfaufg}
% 2016-OS-B1g-m1 (Schwere I, rang 14)
\begin{pfaufg}{nach~P10\\’16~\pkt{}~1g}{5.}{}
Der Scheitelpunkt einer verschobenen Normalparabel ist $S(4\,|\,2)$. Kreuze die Gleichung an, die passt:
\par\smallskip\noindent\pfkreuz{}{$y = (x+4)^2 + 2$}\pfkreuz{}{$y = (x-4)^2 - 2$}\pfkreuz{}{$y = (x-4)^2 + 2$.}
\fusshilfe{\gfuss{5}{(x $-$ 4)² + 2}}
\end{pfaufg}
% 2015-OS-B1i-m1 (Schwere I, rang 15)
\begin{pfaufg}{nach~P10\\’15~\pkt{}~1i}{6.}{}
Die gezeichnete Normalparabel wird um 3 Einheiten nach rechts verschoben. Gib den Scheitelpunkt der verschobenen Parabel an: $S(\ \ |\ \ )$.
\par\smallskip\noindent\gpassend{\begin{ksys}[xmin=-2,xmax=2,ymin=-1,ymax=4,karo=0.5]\parabel{1}{0}{0}{}\end{ksys}}\par
\karo{2}
\fusshilfe{\gfuss{6}{S(3|0)}}
\end{pfaufg}
% 2020-OS-K3d-m1 (Schwere II, rang 16)
\begin{pfaufg}{nach~P10\\’20~\pkt{}~3d}{7.}{}
Die Parabel zu $f$ mit $y = (x+3)^2 - 5$ ist im Koordinatensystem gezeichnet. Spiegelt man den Graphen von $f$ an der y-Achse, entsteht die Parabel $p$. Skizziere $p$ in das Koordinatensystem. Gib eine Gleichung von $p$ an.
\par\smallskip\noindent\gpassend{\begin{ksys}[xmin=-7,xmax=7,ymin=-6,ymax=5,karo=0.5]\parabel{1}{-3}{-5}{}\node[fill=white,inner sep=1pt,font=\small] at (-6.5,3) {$f$};\end{ksys}}\par
\karo{2}
\fusshilfe{\gfuss{7}{(x $-$ 3)² $-$ 5}}
\end{pfaufg}
% 2018-OS-K5c-m1 (Schwere II, rang 18)
\begin{pfaufg}{nach~P10\\’18~\pkt{}~5c}{8.}{}
Im Koordinatensystem ist die Parabel $p$ mit $p(x) = x^2 - 6x + 6$ gezeichnet. Schreib die Gleichung von $p$ in Scheitelpunktform.
\karo{2}
\fusshilfe{\gfuss{8}{(x $-$ 3)² $-$ 3}}
\end{pfaufg}
% 2017-OS-K5e-m1 (Schwere II, rang 19)
\begin{pfaufg}{nach~P10\\’17~\pkt{}~5e}{9.}{}
Die Parabel $p$ hat die Gleichung $y = (x+5)^2 - 3$ und ist im Koordinatensystem gezeichnet. Man verschiebt $p$ um 3 Einheiten nach oben und spiegelt das Ergebnis danach an der x-Achse. So entsteht die Parabel $q$. Gib eine Gleichung von $q$ an.
\par\smallskip\noindent\gpassend{\begin{ksys}[xmin=-9,xmax=1,ymin=-4,ymax=4,karo=0.5]\parabel{1}{-5}{-3}{p}\end{ksys}}\par
\karo{2}
\fusshilfe{\gfuss{9}{$-$(x + 5)²}}
\end{pfaufg}
% 2015-OS-K4b-m1 (Schwere II, rang 20)
\begin{pfaufg}{nach~P10\\’15~\pkt{}~4b}{10.}{}
Lena verlässt ein Flugzeug mit dem Fallschirm. Die ersten 25 Sekunden stürzt sie frei, der Schirm ist noch zu; ihre Höhe sinkt dabei von 2750 m auf 1500 m. Der Graph zeigt ihre Höhe $h$ (in Metern) zur Zeit $t$ (in Sekunden). Welche Gleichung kann den Graphen in den ersten 25 Sekunden beschreiben? Kreuze an und begründe:
\par\smallskip\noindent\gpassend{\begin{ksys}[xmin=0,xmax=300,ymin=0,ymax=3000,xstep=25,ystep=250,karo=0.45,xlabel=$t$ in s,ylabel=$h$ in m]\funktionab{-2*\x^2+2750}{}{0}{25}\funktionab{-6*\x+1650}{}{25}{275}\end{ksys}}\par
\par\smallskip\noindent\pfkreuz{}{$h(t) = 2t^2 + 2750$}\pfkreuz{}{$h(t) = 2t^2 - 2750$}\pfkreuz{}{$h(t) = -2t^2 + 2750$}\pfkreuz{}{$h(t) = -2t^2 - 2750$.}
\fusshilfe{\gfuss{10}{$-$2t² + 2750}}
\end{pfaufg}
% 2014-OS-K7b-m1 (Schwere II, rang 22)
\begin{pfaufg}{nach~P10\\’14~\pkt{}~7b}{11.}{}
Im Koordinatensystem sind die Parabel $p$ mit $p(x) = -x^2$ und die Gerade $g$ mit $g(x) = 3x - 4$ gezeichnet. Eine weitere Gerade $f$ soll die Parabel $p$ nirgends treffen. Gib eine Gleichung für eine solche Gerade an.
\par\smallskip\noindent\gpassend{\begin{ksys}[xmin=-6,xmax=6,ymin=-6,ymax=6,karo=0.5]\parabel{-1}{0}{0}{p}\gerade[2.5]{3}{-4}{g}\end{ksys}}\par
\karo{3}
\fusshilfe{\gfuss{11}{1}}
\end{pfaufg}
% 2018-OS-K5d-m1 (Schwere III, rang 23)
\begin{pfaufg}{nach~P10\\’18~\pkt{}~5d}{12.}{}
Im Koordinatensystem ist die Parabel $p$ mit $p(x) = x^2 - 6x + 6$ gezeichnet. Eine Normalparabel $q$ hat ihren Scheitelpunkt auf der y-Achse und keine Nullstelle. Skizziere eine solche Parabel in das Koordinatensystem und gib ihre Gleichung an. Spiegelt man $q$ an der x-Achse, entsteht $q^{\prime}$. Gib eine Gleichung von $q^{\prime}$ an.
\par\smallskip\noindent\gpassend{\begin{ksys}[xmin=-3,xmax=7,ymin=-4,ymax=7,karo=0.5]\parabel{1}{3}{-3}{p}\end{ksys}}\par
\karo{3}
\fusshilfe{\gfuss{12}{$-$x² $-$ 2}}
\end{pfaufg}
\par\vfill\noindent{\scriptsize\color{mbgrau}Originale: P10 2021 \pkt{} 1d, P10 2020 \pkt{} 3a, P10 2018 \pkt{} 5b, P10 2017 \pkt{} 5c, P10 2016 \pkt{} 1g, P10 2015 \pkt{} 1i, P10 2020 \pkt{} 3d, P10 2018 \pkt{} 5c, P10 2017 \pkt{} 5e, P10 2015 \pkt{} 4b, P10 2014 \pkt{} 7b, P10 2018 \pkt{} 5d\par}
\end{document}